% ============================================================
%  Experiment 1 --- DC Motor / PID Control (Speed)
% ============================================================
\experimentheader{1}{DC Motor Speed Control}{DC Motor / PID Control --- Speed}

\begin{objectives}
\begin{enumerate}
  \item To study the open-loop speed control of a DC motor.
  \item To study closed-loop speed control using P, PD, PI, and PID controllers.
\end{enumerate}
\end{objectives}

\begin{equipment}
\begin{itemize}
  \item Controller kit
  \item Cathode ray oscilloscope
  \item BNC connectors with leads
  \item Multimeter
\end{itemize}
\end{equipment}

\begin{procedure}

\textbf{(a) Open-loop control of DC motor}

\begin{enumerate}
  \item Switch on the supply and measure all constant voltages; calibrate variable voltages in terms of voltage and angle (degrees).
  \item Disconnect all cable wires from the hardware module.
  \item Wire the circuit shown in Figure~\ref{fig:exp1-openloop} for open-loop control.
\end{enumerate}

\begin{boxedfigure}
\centering
\begin{tikzpicture}[node distance=8mm and 10mm]
  \node[io] (vin) {$V_{in}$};
  \node[block, right=of vin] (plant) {Motor\\Plant};
  \node[block, right=of plant] (fvc) {Frequency to\\Voltage Controller};
  \node[block, right=of fvc, minimum width=3cm] (scope) {Oscilloscope};

  \draw[arr] (vin) -- (plant);
  \draw[arr] (plant) -- (fvc) node[midway, above, font=\scriptsize] {Ch-2};
  \draw[arr] (fvc) -- (scope) node[midway, above, font=\scriptsize] {Ch-1};

  % feedback line from plant output up and into scope Ch-1
  \draw[arr] (plant.north) -- ++(0,0.9) -| (scope.north);
\end{tikzpicture}
\captionof{figure}{Open-loop control of motor speed.}
\label{fig:exp1-openloop}
\end{boxedfigure}

\begin{enumerate}[resume]
  \item Connect \textit{Output A} of the motor plant to Channel 2 of the oscilloscope.
  \item Connect \textit{Adjust A} of the motor plant to \textit{Input A} (speed control of the plant) and to Channel 1 of the oscilloscope (only positive voltages are considered).
  \item Turn ON the power switch.
  \item Slowly turn potentiometer-A clockwise ($0$ to $+50$) and anti-clockwise ($0$ to $-50$). Observe and describe the motor operation.
  \item Disconnect \textit{Adjust A} and connect the motor \textit{Input A} to \textit{Adjust B} (only positive voltages are considered).
  \item Vary supply B from $30^\circ$ to $270^\circ$ in steps of $30^\circ$ and study the variation in reference speed (voltage) and output speed (voltage). Plot the error curve.
  \item Comment on your observations and results.
\end{enumerate}

\begin{boxedtable}
\centering
\begin{tabular}{ccccc}
\toprule
Position & $V_{in}$ (Channel I) & $V_{out}$ (Channel II) & Error ($V_{out}-V_{in}$) & Remarks \\
\midrule
$+270$ & & & & \\
$+240$ & & & & \\
$+210$ & & & & \\
$+180$ & & & & \\
$+150$ & & & & \\
$+120$ & & & & \\
$+90$  & & & & \\
$+60$  & & & & \\
$+30$  & & & & \\
\bottomrule
\end{tabular}
\captionof{table}{Open-loop speed response.}
\end{boxedtable}

\textbf{(b) Closed-loop control of speed of DC motor}

\textit{(i) Using P-controller with feedback through optical encoder}\vspace{2pt}

\begin{enumerate}
  \item Select the P-controller by setting PROP on the controller (DIP switch 1) to ON.
  \item Connect \textit{Adjust B} of the plant to the input of the PID controller.
  \item Connect \textit{Output A} of the plant to the controller feedback.
  \item Connect the controller output to \textit{Input A} of the plant.
  \item Turn potentiometer-B of the motor plant to a middle position (e.g.\ $+180$).
  \item Adjust trimmer $R_{p2}$ (varies $K_p$) to obtain an optimum response (minimum steady-state error).
  \item Turn potentiometer-B to vary the speed of the motor. Record the input and output voltages for various positions in tabular form.
\end{enumerate}

\textit{(ii) Speed control with PI-controller}\vspace{2pt}

\begin{enumerate}
  \item On the PID controller card, set trimmers $R_p$ and $R_d$ to minimum (counter-clockwise) while $R_i$ is set to maximum (clockwise).
  \item Select the PI-controller by setting PROP and INTG to ON.
  \item Slowly vary trimmer $R_i$ to obtain an optimum output response (minimum steady-state error). If gain is insufficient (or overshoots), adjust $R_{p2}$.
  \item Turn potentiometer-B to vary the speed of the motor. Record the input and output voltages for various positions (Table~\ref{tab:exp1-closed}). Repeat step 3 if necessary to fine-tune the response.
  \item Compare the experimental results for the P-controller and the PI-controller.
\end{enumerate}

\textit{(iii) Speed control with PD-controller}\vspace{2pt}

\begin{enumerate}
  \item Switch on the supply and measure all constant voltages; calibrate variable voltages in terms of voltage and angle.
  \item On the PID controller card, set trimmers $R_p$ and $R_d$ to minimum (counter-clockwise) while $R_i$ is set to maximum (clockwise).
  \item Select the PD-controller by setting PROP and DIFF to ON; all other switches OFF.
  \item Vary supply B from $30^\circ$ to $270^\circ$ in steps of $30^\circ$ and study the variation in reference speed ($V_{in}$) and output speed ($V_o$). Plot the error curve.
  \item Study the effect of varying $R_d$ on the output.
\end{enumerate}

\textit{(iv) Speed control with PID-controller}\vspace{2pt}

\begin{enumerate}
  \item Adjust potentiometer-B to $2\,\text{V}$ (around $+180^\circ$).
  \item Set trimmer $R_{p2}$ to minimum to reduce the damping ratio.
  \item Set trimmer $R_i$ to minimum. The output should start to vibrate.
  \item Slightly increase $R_i$ until the system starts oscillating between two points (pure oscillation).
  \item Slightly increase $R_{p2}$ to bring the system to a damped condition; it should oscillate for a few cycles.
  \item Apply a small force to hold the motor (the oscilloscope will show the motor output drop to 0\,V); observe and describe the response when the force is released.
  \item Select the PID-controller by setting PROP, INTG, and DIFF to ON.
  \item Slowly vary $R_d$ to obtain the optimum response.
  \item Compare its operation with the P, PI, and PD controllers under disturbance.
  \item As before, take observations by varying the applied voltage.
\end{enumerate}

\begin{boxedtable}
\centering
\begin{tabular}{ccccc}
\toprule
Position & $V_{in}$ (Channel I) & $V_{out}$ (Channel II) & Error ($V_{out}-V_{in}$) & Remarks \\
\midrule
$+270$ & & & & \\
$+240$ & & & & \\
$+210$ & & & & \\
$+180$ & & & & \\
$+150$ & & & & \\
$+120$ & & & & \\
$+90$  & & & & \\
$+60$  & & & & \\
$+30$  & & & & \\
\bottomrule
\end{tabular}
\captionof{table}{Closed-loop speed response.}
\label{tab:exp1-closed}
\end{boxedtable}

\end{procedure}

\begin{reportq}
\begin{enumerate}
  \item Compare the performance of P, PI, PD, and PID controllers.
  \item What will happen if all types of controllers are connected in series?
  \item What are the different types of encoders? Explain the working principle of an optical encoder with the help of suitable diagrams.
\end{enumerate}
\end{reportq}
